\chapter{Appendix C: Model Geometry, Bracket Frames, and Compliance Parameters}\label{app:model_geometry}

This appendix documents the numeric geometry and compliance parameters behind the torque models of Chapters~\ref{ch:methods}--\ref{ch:results} as they exist in the project's MATLAB code, with file names given so that any value can be traced to the code that produced it. Route and bracket coordinates are coded in meters and reported here in millimeters; stiffnesses are coded and reported in newtons per meter. The bracket reference points are interpretation choices for the 3D-printed (Onyx FDM) attachment hardware, set from CAD by the author; moving a reference point changes the identified stiffness values, so the points are recorded here alongside the values they produced.

\section{Static BPA Pathway Data}

The static BPA model routes each actuator as a straight-segment polyline through via points stored in a per-test \texttt{Location} array of dimensions $M\times3\times N$, where $N$ is the number of joint frames in the motion sweep and $M$ the number of via points; the first row in the distal link frame is marked by a cross-point index. There are no binormal or surface-frame quantities: at each frame the force direction is the normalized vector from the cross-point row to the previous row transformed into the crossing frame, $\hat{\mathbf{u}} \propto \mathbf{R}^{-1}\,\mathbf{p}_{C-1} - \mathbf{p}_{C}$ [Provenance: \texttt{Testing\_Data/2022\_02\_Festo/minimizeFlxPin2brk.m} lines 429--440; \texttt{minimizeFlxPin.m} lines 381--394; \texttt{minimizeExtX3.m} lines 676--687; BPA struct fields, \texttt{minimizeFlxPin.m} lines 28--35]. The pinned-knee test routes are archived as \texttt{FlxPinBPASet.mat} (5 flexor tests, $M=2$, knee angle $-120^{\circ}$ to $+10^{\circ}$) and \texttt{ExtPinBPASet.mat} (9 extensor tests, $M=7$, knee angle $-125^{\circ}$ to $+35^{\circ}$).

Table~\ref{tab:app_flxpin_tests} lists the five pinned-knee flexor tests; their routes are two points, femur-frame origin $\mathbf{p}_{1} = [-75.0,\,100.0,\,32.8]$~mm and tibia-frame insertion $\mathbf{p}_{2} = [-50.1,\,-45.0,\,32.3]$~mm (test 1; \texttt{FlxPinBPASet.mat}, kf(1), frame 1). Table~\ref{tab:app_extpin_route} lists the extensor pinned-knee route of test 1 at $\theta_{k}=-125^{\circ}$; rows 1--6 are femur-frame and row 7 is tibia-frame, and as the knee extends, rows 2--6 release onto $\mathbf{p}_{1}$ (wrap released at $+35^{\circ}$). Table~\ref{tab:app_extpin_tests} lists the per-test constants of the nine extensor tests; the validation pool used for the correction terms of Section~\ref{sec:improved_model} is tests \{1, 2, 5, 6, 7, 8\} [Provenance: \texttt{crossPredictFlx.m} line 46].

\begin{sidewaystable}
\caption{Pinned-knee flexor test BPA constants (\texttt{FlxPinBPASet.mat}; test labels from \texttt{minimizeFlxPin10mm\_2brk.m} line 42). All tests use $\phi$\qty{10}{\mm} BPAs; \texttt{fitn} is the end-cap fitting length.}
\label{tab:app_flxpin_tests}
\begin{tabular}{@{}llllll@{}}
\toprule
\textbf{Test} & $\boldsymbol{l_{0}}$ (m) & Tendon (m) & $P$ (kPa) & $F_{m}$ (N) & \textbf{Notes} \\
\midrule
48\,cm        & 0.485 & 0     & 605.7 & 443.4 & \\
46\,cm        & 0.457 & 0     & 614.2 &       & \\
47\,cm        & 0.479 & 0     & 623.2 &       & \\
40\,cm-tendon & 0.406 & 0.035 & 623.4 &       & \\
41\,cm        & 0.410 & 0     & 615.4 &       & \\
\bottomrule
\end{tabular}
\end{sidewaystable}

\begin{sidewaystable}
\caption{Biomimetic knee, extensor 20\,mm BPA route seed (femur frame: p1--p5; tibia frame: p6--p9; \texttt{buildKneeExtContext20mm.m} lines 349--384, Ben-set SolidWorks points). Release schedule of the verified elimination: p5 at $-91.1^{\circ}$, p4 at $-50.4^{\circ}$, p6 at $-47.8^{\circ}$, p7 at $-20.2^{\circ}$, p3 at $+6.1^{\circ}$.}
\label{tab:app_ext_seed}
\begin{tabular}{@{}lrrrll@{}}
\toprule
\textbf{Point} & $x$ (mm) & $y$ (mm) & $z$ (mm) & \textbf{Frame} & \textbf{Role} \\
\midrule
p1 & 40.00 & 35.00  & 0 & femur & BPA origin (attachment) \\
p2 & 83.90 & $-274.76$ & 0 & femur & BPA contacts mounting base \\
p3 & 66.34 & $-427.24$ & 0 & femur & femoral contact \\
p4 & 51.52 & $-446.74$ & 0 & femur & femoral condyle contact \\
p5 & 38.26 & $-453.01$ & 0 & femur & femoral condyle contact \\
p6 & 69.85 & 34.33  & 0 & tibia & tibial contact (upper) \\
p7 & 73.50 & 19.13  & 0 & tibia & tibial contact (mid) \\
p8 & 71.98 & $-12.23$ & 0 & tibia & tibial contact (lower) \\
p9 & 34.06 & $-66.87$ & 0 & tibia & patellar ligament ring (distal, pEnd) \\
\bottomrule
\end{tabular}
\end{sidewaystable}

\begin{table}[htbp]
\caption{Pinned-knee extensor BPA route of test 1 at $\theta_{k}=-125^{\circ}$ (\texttt{ExtPinBPASet.mat}, ke(1), frame 1; mm).}
\label{tab:app_extpin_route}
\begin{tabular}{@{}lrrr@{}}
\toprule
\textbf{Point} & $x$ (mm) & $y$ (mm) & $z$ (mm) \\
\midrule
p1 & 30.00  & $-50.00$ & 0 \\
p2 & 65.00  & $-350.00$ & 0 \\
p3 & 42.40  & $-408.50$ & 0 \\
p4 & 28.20  & $-421.60$ & 0 \\
p5 & 10.50  & $-425.20$ & 0 \\
p6 & $-8.60$ & $-422.80$ & 0 \\
p7 & 42.50  & $-75.90$ & 0 \\
\bottomrule
\end{tabular}
\end{table}

\begin{table}[htbp]
\caption{Pinned-knee extensor test BPA constants (\texttt{ExtPinBPASet.mat}, ke(2)--ke(9); test 1: $l_0 = 0.400$~m, tendon 0, $P = 624.9$~kPa, $F_m = 436.3$~N). All tests use $\phi$\qty{10}{\mm} BPAs.}
\label{tab:app_extpin_tests}
\begin{tabular}{@{}llll@{}}
\toprule
\textbf{Test} & $\boldsymbol{l_{0}}$ (m) & Tendon (m) & $P$ (kPa) \\
\midrule
40\,cm-tendon & 0.400 & 0.040 & 605.2 \\
42\,cm        & 0.420 & 0     & $\sim$620 \\
42\,cm-tendon & 0.420 & —     & $\sim$620 \\
43\,cm        & 0.430 & 0     & $\sim$620 \\
43\,cm-tendon & 0.430 & —     & $\sim$620 \\
46\,cm        & 0.460 & 0     & $\sim$620 \\
47\,cm        & 0.470 & 0     & $\sim$620 \\
48\,cm        & 0.480 & 0     & $\sim$620 \\
\bottomrule
\end{tabular}
\end{table}

The biomimetic flexor route (\texttt{buildKneeFlexorContext20mm.m}) is a three-point path --- femur-frame origin $\mathbf{p}_{1}$, a moving wrap point $\mathbf{p}_{\text{wrap}}$ on the tibia, and tibia-frame insertion $\mathbf{p}_{\text{end}}$ --- with the wrap point held on a standoff surface of radius \texttt{bpaRb} about the joint and its $z$ placed on the projected $\mathbf{p}_{1}$--$\mathbf{p}_{\text{end}}$ line; the optimizer seed is $\mathbf{p}_{1,0} = [-50.0,\,35.0,\,50.0]$~mm and $\mathbf{p}_{2,0} = [-12.24,\,-8.87,\,27.87]$~mm with $l_{0} = 415$~mm and 25~mm tendon [Provenance: \texttt{buildKneeFlexorContext20mm.m} lines 96--119, 179--188; \texttt{buildKneeFlexorRoute20mm.m} lines 99--186]. The biomimetic extensor seed route is given in Table~\ref{tab:app_ext_seed}, including the release schedule of the route-elimination rule.

\section{Joint Transformation Matrices}

The knee transform is planar: every joint frame is a rotation about $z$ by the knee angle $\varphi$ composed with a translation fitted to robot CAD tables of the migrating instantaneous center of rotation. With $\mathbf{R} = \mathbf{R}_{z}(\varphi)$, the robot (femur$\to$ICR) transform is $\mathbf{T}_{\text{Pam}} = [\mathbf{R},\; [f_{3}(\varphi),\,f_{4}(\varphi),\,0]^{\top};\; \mathbf{0}^{\top}\;1]$ with $f_{3}, f_{4}$ cubic-spline fits (ICR $x$: 23.30 to 10.47~mm; $y$: $-416.65$ to $-395.66$~mm over $\varphi = 0.17$ to $-2.62$~rad), and the ICR$\to$tibia transform is the pure translation $[f_{13}(\varphi),\,f_{14}(\varphi),\,0]^{\top}$ ($x$: 29.66 to $-7.58$~mm; $y$: 25.97 to 22.33~mm), composed as $\mathbf{T}_{t1\!\to\!f} = \mathbf{T}_{t1\!\to\!\text{ICR}}\;\mathbf{T}_{\text{Pam}}^{-1}$ over 100 frames spanning $\varphi \in [-120^{\circ}, +10^{\circ}]$ [Provenance: \texttt{buildKneeFlexorContext20mm.m} lines 34--80; \texttt{buildKneeExtContext20mm.m} lines 30--78; \texttt{Functions/ModernRobotics/RpToTrans.m}]. In the archived test sets the per-frame transform maps tibia to femur/ICR; at $\theta_{k}=-120^{\circ}$ (test kf1),
\begin{equation}
    \mathbf{T}_{k} \;=\;
    \begin{bmatrix}
    -0.5000 & 0.8660 & 0 & 0.0045\\
    -0.8660 & -0.5000 & 0 & -0.3958\\
    0 & 0 & 1 & 0\\
    0 & 0 & 0 & 1
    \end{bmatrix},
    \label{eq:app_Tk}
\end{equation}
i.e., $\mathbf{R}_{z}(-120^{\circ})$ with constant translation $[4.5,\,-395.8,\,0]$~mm, rotating to $\mathbf{R}_{z}(+10^{\circ})$ at full extension (\texttt{FlxPinBPASet.mat}, kf(1)).

\section{Bracket Reference Points and Frames}

Table~\ref{tab:app_bracket_points} lists the bracket reference points. Each compliant bracket carries an instantaneous body-fixed frame anchored at its reference point, constructed by two rotations (yaw about $z$, then pitch about the yawed $y$): with $\mathbf{p}_{b}$ the vector from the reference point to the bracket-to-BPA attachment target, $\theta_{br} = \operatorname{atan2}(p_{by}, p_{bx})$, $\mathbf{R}_{kbr} = \mathbf{R}_{z}(\theta_{br})\,\mathbf{R}_{y}(\phi_{br})^{\top}$ with $\phi_{br} = \operatorname{atan2}(p'_{bz}, p'_{bx})$ measured on the yawed vector $\mathbf{p}'_{b} = \mathbf{R}_{z}^{\top}\mathbf{p}_{b}$, and $\mathbf{T}_{kbr} = [\mathbf{R}_{kbr},\, \mathbf{p}_{br};\, \mathbf{0}\;1]$; the frame's $x$ axis therefore aims along the nominal load path. The angles are computed per configuration from the bracket geometry, not fixed [Provenance: \texttt{minimizeFlxPin2brk.m} lines 175--203; \texttt{minimizeExtX3.m} lines 463--474]. Force is transformed into the bracket frame by $\mathbf{RowVecTrans}(\mathbf{T}_{kbr}^{-1}, \mathbf{F}_{k})$ and deformed points are mapped back by $\mathbf{RowVecTrans}(\mathbf{T}_{kbr}, \cdot)$. The single-rotation (yaw-only) variant used in the earlier one-transform fit survives in \texttt{minimizeFlxPinX3.m} line 205.

\begin{table}[htbp]
\caption{Bracket reference points (mm, as coded; divided by 1000 in code). All are Ben-set CAD interpretation points; frame denotes the frame the vector is expressed in.}
\label{tab:app_bracket_points}
\begin{tabular}{@{}lrrll@{}}
\toprule
\textbf{Point} & $x$ & $y$ & $z$ & \textbf{Configuration / meaning} \\
\midrule
Pbri & $-27.50$ & $-107.81$ & $-0.54$ & pinned-knee flexor, insertion bracket (knee-ICR frame) \\
Pbr2 & $-52.61$ & 0 & 75.06 & pinned-knee flexor, origin-side 2nd bracket (hip frame) \\
Pbr  & $-3.84$  & $-46.44$ & 62.50 & pinned-knee extensor, rib midpoint, medial side (hip frame) \\
Pbr  & 8.38 & 20.75 & 25.10 & biomimetic extensor 10\,mm, cantilever centroid \\
Pbr  & 9.48 & $-33.38$ & 33.90 & biomimetic flexor 20\,mm, bolt-pattern centroid \\
Pbr  & $-19.00$ & 22.00 & 27.60 & biomimetic flexor 10\,mm, cantilever centroid \\
\bottomrule
\end{tabular}
\end{table}

\section{Stiffness Arrays and Compliance}

Every configuration models bracket compliance as a diagonal Hooke-law spring in the bracket frame, $\mathbf{K}_{br} = \operatorname{diag}(k_{1}, k_{2}, k_{3})$ newtons per meter, with the scalar compliance projected onto the force-path unit vector, $c_{b} = \hat{\mathbf{u}}\,\mathbf{K}_{br}^{-1}\hat{\mathbf{u}}^{\top}$, exactly as in Equations~(\ref{eq:bktstiffness})--(\ref{eq:keff}). The per-configuration stiffness arrays and their series combinations are given in Table~\ref{tab:app_stiffness_arrays}; when an artificial tendon is present its axial spring ($\times2$ cables, effective steel area $1.51\times10^{-6}$~m$^{2}$, $E = 193$~GPa) adds in series, $c_{\text{eff}} = c_{b} + c_{\text{ten}}$ ($+\,c_{b2}$ for the second bracket), and the force balance $F_{\max}F^{*}(\varepsilon^{*}(\delta), P^{*}) = k_{\text{eff}}\,\delta$ is solved per frame by \texttt{fzero} on the contraction interval [Provenance: \texttt{minimizeFlxPin2brk.m} lines 284--339, 489--499; \texttt{minimizeFlxPin.m} lines 266--275; \texttt{minimizeExtX3.m} lines 562--570; \texttt{minimizeExt.m} line 352; \texttt{minimizeFlx.m} line 314].

\begin{table}[htbp]
\caption{Bracket stiffness arrays by configuration ($K = [k_{1}, k_{2}, k_{3}]$ in bracket-frame axes, N/m; optimizer solves $X_{1}, X_{2}$ assigned as shown).}
\label{tab:app_stiffness_arrays}
\begin{tabular}{@{}lll@{}}
\toprule
\textbf{Configuration} & $\boldsymbol{K}$ array & \textbf{Source} \\
\midrule
Pinned flexor, single bracket       & $[X_{1},\,X_{2},\,X_{2}]$ & \texttt{minimizeFlxPin.m} line 266 \\
Pinned flexor, insertion bracket    & $[X_{1},\,X_{2},\,X_{1}]$ & \texttt{minimizeFlxPin2brk.m} line 284 \\
Pinned flexor, origin bracket       & $[X_{2},\,X_{1},\,X_{2}]$ & \texttt{minimizeFlxPin2brk.m} line 298 \\
Pinned extensor (with $X_{3}$)      & $[X_{2},\,X_{1},\,X_{2}]$ & \texttt{minimizeExtX3.m} line 562 \\
Biomimetic extensor                 & $[X_{1},\,X_{2},\,X_{2}]$ & \texttt{minimizeExt.m} line 352 \\
Biomimetic flexor                   & $[X_{1},\,X_{2},\,X_{2}]$ & \texttt{minimizeFlx.m} line 314 \\
\bottomrule
\end{tabular}
\end{table}

The wrapping-loss term $X_{3}$ (Equations~\ref{eq:deltaL}--\ref{eq:x3}) is evaluated as $X_{3}\,R\,\lvert\Delta\theta\rvert\,\zeta^{2}$ per contact arc. On the pinned extensor the two physical contact arcs use $R_{1} = 12$~mm ($27^{\circ}$ to $-23^{\circ}$) and $R_{2} = 40$~mm, with arc-length normalizers $27.75/(ang_{1} - ang_{2})$ and $92.46/(ang_{2}+120^{\circ})$; on the biomimetic configurations the bend measure is $w_{\text{rap}}\lvert\theta_{\text{turn}}\rvert$ with $w_{\text{rap}} = 25$~mm [Provenance: \texttt{minimizeExtX3.m} lines 333--359; \texttt{buildKneeFlexorRoute20mm.m} line 268; \texttt{MonoPamDataExplicit\_balanceX3.m} lines 489--494]. The optimizer parameterization is $[100\,X_{0}\ (\text{cm}),\ \log_{10}X_{1},\ \log_{10}X_{2},\ \ldots]$ with bounds $X_{0}\in[0, 20]$~mm and $X_{1}, X_{2}\in[5\times10^{3}, 5\times10^{7}]$~N/m for the flexor; the extensor refit locks $(X_{1}, X_{2})$ to the flexor values and solves $X_{0}\in[-20, 0]$~mm and $X_{3}\in[0, 1]$, enforcing a single stiffness pair across configurations [Provenance: \texttt{minimizeFlxPin10mm\_2brk.m} lines 62--63, 402--404; \texttt{minimizeExt10mmX3.m} lines 46--47; \texttt{sweepExtX3.m} lines 59--60].

\section{Adopted Correction-Term Values}

Table~\ref{tab:app_xi_values} records the correction-term values adopted by the archived code at the time of writing. The pinned-knee fits of record for the published comparison are those of \texttt{minimizeExtPin10\_results\_20260819\_2transforms\_Z2.mat}; the current redesign operating point (context builders \texttt{buildKneeFlexorContext20mm.m} lines 150--165 and \texttt{buildKneeExtContext20mm.m} lines 150--160) pairs the two-bracket flexor fit with an extensor refit in which $(X_{1}, X_{2})$ are locked to the flexor values. Regenerating the results files updates this table.

\begin{table}[htbp]
\caption{Correction-term values of record and current operating point. $(X_{1}, X_{2})$ are shared across flexor and extensor in the current operating point.}
\label{tab:app_xi_values}
\begin{tabular}{@{}llll@{}}
\toprule
\textbf{Term} & \textbf{Pinned fit of record} & \textbf{Current operating point} & \textbf{Source (.mat)} \\
\midrule
Flexor $X_{0}$ (m)   & ---                & $1.647\times10^{-6}$ & crossPredict\_20260907 \\
Flexor $X_{1}$ (N/m) & ---                & $2.107\times10^{6}$  & crossPredict\_20260907 \\
Flexor $X_{2}$ (N/m) & ---                & $1.065\times10^{4}$  & crossPredict\_20260907 \\
Extensor $X_{0}$ (m) & $-1.220\times10^{-2}$ & $-2.364\times10^{-3}$ & both \\
Extensor $X_{1}$ (N/m) & $1.508\times10^{4}$ & $= $ flexor        & both \\
Extensor $X_{2}$ (N/m) & $1.218\times10^{4}$ & $= $ flexor        & both \\
$X_{3}$ (--)         & 0.2937             & 0.4637 (refit)       & both \\
\bottomrule
\end{tabular}
\end{table}
